\section{Going faster}\label{part:faster}

Everything up to here is enough to explain the \emph{correctness} of the core:
in an unlimited exact run, the solver of \crefrange{part:foundations}{part:together}
finds optima and proves them. This section is about being \emph{fast}. Each
technique that follows is an accelerator layered onto that same core: the LP
relaxation, cutting planes, restarts, and the parallel portfolio. In the idealized complete-search sense, these do not change the
feasible set or the optimum; in practical time-limited runs, they can of course
change whether CP-SAT proves optimality, returns only an incumbent, or times out
without a useful solution.

\begin{platypusbook}[htb]
The accelerators in this section all run \emph{during} search. CP-SAT also spends
a considerable effort \emph{before} search begins, on \textbf{presolve}: a
preprocessing pass that simplifies and canonicalizes the model, removes
redundant variables and constraints, tightens bounds, and expands high-level
constraints into the primitives the search core understands. None of it is
exotic, but it is fundamental to performance, and a model often shrinks
dramatically before a single decision is made. We leave presolve out of scope
here; the general ideas, drawn from the SAT and MIP traditions CP-SAT inherits,
are covered in lectures on SAT and MaxSAT
preprocessing~\cite{preprocsat,preprocmaxsat} and in Achterberg et al.'s survey
of MIP presolve reductions~\cite{achterberg2016presolve}.
\end{platypusbook}

\subsection{The LP relaxation: dual simplex as a propagator}\label{sec:lp}
% Sources (VERIFIED):
%   ortools/sat/linear_programming_constraint.h -- class LinearProgrammingConstraint
%     : public PropagatorInterface, ReversibleInterface; Propagate()/IncrementalPropagate();
%     GetBasisState()/LoadBasisState() (warm start); header comment "uses glop's
%     revised simplex ... Reduced Cost Strengthening".
%   ortools/sat/linear_programming_constraint.cc --
%     simplex_params_.set_use_dual_simplex(true)  (uses DUAL simplex);
%     watcher_->Register(this) + RegisterWith at propagation priority 2;
%     state_=simplex_.GetState() / simplex_.LoadStateForNextSolve(state_) (warm start);
%     AnalyzeLp(): Enqueue(GreaterOrEqual(objective_cp_, ceil(lp_obj))) (dual bound);
%     ReducedCostStrengtheningDeductions() (reduced-cost fixing);
%     PropagateExactLpReason(): scale negated dual values to integers (power-of-2,
%     divide by GCD) -> exact integer linear combination = Farkas-style reason;
%     root-level dtime throttling + max_number_of_iterations limits.
%   ortools/glop/parameters.proto -- use_dual_simplex (Glop revised simplex).
%   ortools/sat/sat_parameters.proto -- linearization_level (default 1; range 0..2;
%     no_lp=0, default_lp=1, max_lp=2 in cp_model_search.cc). NOTE: the developer
%     slides show a level 3, but the released code tops out at 2.
%     use_exact_lp_reason (default true).
%   linear_programming_constraint.cc -- AnalyzeLp() at DUAL_UNBOUNDED calls
%     PropagateExactDualRay() (GetDualRay) -> conflict; the third LP deduction.

Pure constraint propagation reasons one constraint at a time. It is weak exactly
where linear programming is strong: at providing a \emph{global} numeric view of
all linear constraints at once, a sharp lower bound on the objective, and
reduced-cost arguments. CP-SAT therefore carries a continuous \textbf{LP
relaxation} of part of the model alongside the search and uses it as one more
propagator. This is the bridge to the MIP world, and it is where the dual
simplex enters. The developers' own account of these MIP aspects is a useful
companion here~\cite{didier2025mipaspects}.

\textbf{It is just a propagator.}
\srcla{linear\_programming\_constraint.h}{137}{LinearProgrammingConstraint}
implements the same \srcla{integer.h}{1276}{PropagatorInterface} as the CP
propagators and registers with the
\srcla{integer.h}{1324}{GenericLiteralWatcher}. It registers at a high-cost
priority, so the cheapest-first rule of \cref{sec:propstep} runs it only
after unit propagation and the light CP propagators have reached their fixed
point. It watches integer-variable bounds (Booleans included, through their
$\{0,1\}$ integer views) and the objective bound, wakes
incrementally when those bounds change, and, like any propagator, feeds its
deductions back as integer-bound tightenings on the trail. Nothing in the search
core changes; the LP is simply an unusually powerful propagator.

\textbf{What it relaxes.} CP-SAT loads the model's linear constraints into a
continuous LP over the same variables, drops integrality, and uses each
variable's current $[\text{lb}, \text{ub}]$ as its column bounds. How much of the
model enters the LP is governed by \code{linearization\_level}: level 0 means no
LP at all; level 1 adds the linear constraints together with the integer and
at-most-one encodings; and level 2 linearizes more aggressively still, pulling in
additional Boolean and reified structure and enabling the cut generators of
\cref{sec:cuts} where they apply. The proto documents only this coarse
0/1/2 progression; the precise set of relaxations admitted at each level is an
implementation detail and shifts between releases. This single knob is what
distinguishes the \code{no\_lp} (0), \code{default\_lp} (1), and \code{max\_lp}
(2) workers of the portfolio (\cref{sec:portfolio}).

\textbf{Why the dual simplex.} This is the crux. Branching and propagation change
only variable \emph{bounds}, the column bounds of the LP; they do not change the
constraint matrix or the objective. After such a bound change, the previous
optimal basis often remains dual-feasible, because dual feasibility depends on
the matrix and objective, while primal feasibility may be broken by the new
bounds. This is the textbook situation for the \textbf{dual simplex}: starting
from the saved basis, it repairs primal feasibility, typically in a small number
of pivots, instead of re-solving from scratch. CP-SAT sets
\code{use\_dual\_simplex(true)} and saves and reloads Glop's basis state across
calls (\srcla{linear\_programming\_constraint.cc}{964}{GetState} / \srcla{linear\_programming\_constraint.h}{269}{LoadStateForNextSolve}). This warm start is what
makes it affordable for an LP to sit inside a CDCL search.

\textbf{What it gives back.} The LP returns three kinds of deduction, each handed
to the engine as ordinary integer-bound facts, or as a conflict, on the trail:

\begin{itemize}
  \item \textbf{A conflict, when the LP is infeasible.} If the relaxation has no
  feasible point, CP-SAT can use a \textbf{dual ray}, a Farkas certificate, to
  form a linear combination proving that the current bounds are jointly
  impossible. That proof is reported as a learned conflict like any other
  (\srcla{linear\_programming\_constraint.h}{329}{PropagateExactDualRay}).
  \item \textbf{A dual bound on the objective.} When the LP is optimal, its value
  is a valid lower bound, for minimization, on the relaxation and therefore on
  the integer problem. The propagator pushes
  $\text{objective\_var} \ge \lceil \text{lp\_objective} \rceil$, directly
  strengthening the bound that drives the optimality proof of
  \cref{sec:opt}.
  \item \textbf{Reduced-cost strengthening.} Given the incumbent cutoff, i.e.,
  the objective's current upper bound, each variable's LP reduced cost says how
  much moving it away from its LP value would worsen the objective. If that
  movement would breach the cutoff, the variable's domain can be tightened. This
  is classic MIP reduced-cost fixing, returned as ordinary integer-bound
  deductions.
\end{itemize}

\textbf{Soundness despite floating point.} The simplex runs in floating point,
which a solver that must be \emph{exact} and must produce \emph{learnable
reasons} cannot trust directly. CP-SAT therefore does not blindly believe the
LP's numerical values; it treats them as proposals. For LP-derived deductions, it
scales the relevant dual multipliers to integers, using a power-of-two scaling
and then dividing by the gcd, and forms an exact integer linear combination of
the constraints. The result is a Farkas-style certificate that yields a provably
valid bound together with an exact reason expressed in integer bound literals.
That certified reason is what is enqueued and, if a conflict follows, learned.
With \code{use\_exact\_lp\_reason} on by default, the \emph{deductions are exact
even though the LP that suggested them was approximate}. This is what keeps the
LP compatible with lazy clause generation: \emph{the LP proposes, exact integer
arithmetic certifies, and the certified reason is there for CDCL when the
deduction enters conflict analysis.}

\textbf{Cost and throttling.} A simplex re-solve is far heavier than an ordinary
propagator, so CP-SAT throttles the LP beyond merely assigning it a low priority.
At the root, it may be skipped unless enough deterministic time has elapsed since
the previous solve; pivot counts are capped per call; and the LP can be disabled
outright during heavy phases such as probing. The portfolio hedges this
trade-off directly: some workers pay nothing for the LP (\code{no\_lp}), whereas
others pay for a maximal relaxation (\code{max\_lp}), so the ensemble covers both
sides of the bet on whether the linear structure is worth the cost.

In short, the LP injects the kind of reasoning that constraint propagation cannot
provide on its own: a global, dual-certified numeric bound. It does so without
disturbing the engine. It is a propagator, woken by bound changes, that returns
exactly justified bound tightenings, made affordable by dual-simplex warm starts.

% Source (VERIFIED): "Lp stats" table header {"Lp stats","Component","Iterations",
%   "AddedCuts","OPTIMAL","DUAL_F.","DUAL_U."} at ortools/sat/stat_tables.cc:49-51;
%   Iterations = lp->total_num_simplex_iterations() (stat_tables.cc:239).
\begin{platypuslog}[htb]
Every LP-bearing worker adds a row to the \textsf{Lp stats} table:
\code{Iterations} is its cumulative simplex pivot count, held low precisely by
the dual-simplex warm starts, and the \code{OPTIMAL}/\code{DUAL\_F.}/\code{DUAL\_U.}
columns count solve outcomes. The \code{DUAL\_U.} column is the dual-unbounded
case, corresponding to the infeasible relaxation that produces the dual-ray
conflict described in this section.
\end{platypuslog}
\subsection{Cutting planes: tightening the relaxation}\label{sec:cuts}
% Sources (VERIFIED):
%   ortools/sat/linear_programming_constraint.cc -- AddCGCuts() (1607, Chvatal-Gomory,
%     multipliers from B^-1.e_j), AddMirCuts() (1804), AddZeroHalfCuts() (2071);
%     generic cuts generated mainly at level zero (root); cut loop ~2215-2275.
%   ortools/sat/cuts.{h,cc} -- IntegerRoundingCutHelper (cuts.cc:1002), CoverCutHelper
%     /knapsack (1534, Letchford-Souli lifting), ImpliedBoundsProcessor (2338),
%     CreateCliqueCutGenerator (2971), BoolRLTCutHelper (cuts.h:637). RHS kept as
%     absl::int128, coefficients IntegerValue -> cuts are EXACT (integer arithmetic).
%   ortools/sat/scheduling_cuts.{h,cc} -- energetic (CreateCumulative/NoOverlapEnergyCutGenerator),
%     completion-time / Smith rule (CreateCumulative/NoOverlapCompletionTimeCutGenerator),
%     time-table, precedence cuts. ortools/sat/diffn_cuts.h -- 2d variants.
%   ortools/sat/routing_cuts.{h,cc} -- subtour/strongly-connected (3196), CVRP (3214),
%     flow cuts; registered in linear_relaxation.cc.
%   ortools/sat/linear_constraint_manager.cc -- AddCut() (422): efficacy threshold +
%     orthogonality filtering; TopN cut pool, number of active cuts capped.

The LP of \cref{sec:lp} is only as strong as the constraints in it, and
the bare relaxation of an integer model is often loose.
% Source (VERIFIED): LowerThanCoeffStrengthening(), ortools/sat/cp_model_presolve.cc:4570
%   -- "coeff * (X - LB + LB) -> rhs * (X - LB) + coeff * LB"; clamps an over-large
%   coefficient down to the rhs (comment ~4759: "3x + 5y <= 6, the coefficient 5 can
%   be changed to 6"); also rewrites through an enforcement literal when implied.
Presolve already buys some tightening for free: \emph{coefficient strengthening}
rewrites a coefficient that the right-hand side leaves partly slack into a
tighter one, without changing the set of integer solutions. The source's own
example is $3x + 5y \le 6$ with $x, y \in \{0,1\}$: the coefficient $5$ can be
raised to $6$, since $y = 1$ already forces $x = 0$ either way, so no integer
point is lost while the linear relaxation becomes sharper at no extra row. True
to the engine's preference for Boolean structure, CP-SAT can even re-express
such a constraint through an \emph{enforcement literal} when one is implied.

\textbf{Cutting planes} sharpen the relaxation further. They are extra linear
inequalities that every \emph{integer} solution satisfies but that the current
fractional LP optimum violates. Adding them cuts off the fractional point without
removing any feasible integer point, so the next re-solve can return a tighter
dual bound. This is the MIP half of CP-SAT: cuts are generated into the same
shared LP and gated by the \code{linearization\_level} of \cref{sec:lp}.

\textbf{Generic MIP cuts.} From the LP basis, CP-SAT separates standard cut
families: \textbf{Chvátal--Gomory} cuts, with multipliers read from $B^{-1}$;
\textbf{MIR} cuts, mixed-integer rounding cuts obtained by combining a few tight
rows; \textbf{zero-half} cuts, based on a mod-2 heuristic; and
\textbf{knapsack/cover} cuts with lifting. It also derives implied-bound and
clique cuts from the Boolean structure.

% Source (VERIFIED): the unifying super-additive recipe is GetSuperAdditiveRoundingFunction()
%   in ortools/sat/cuts.cc:709-792 -- MIR function f(x) = size*floor(t*x/divisor)
%   + max(0, (t*x mod divisor) - rhs_remainder); declared cuts.h:310-343 with the
%   comment "f is super-additive: f(a) + f(b) <= f(a + b)". CoverCutHelper::
%   TrySimpleKnapsack (cuts.cc:1534) complements the cover then applies it with
%   divisor = best_coeff (cuts.cc:1629). Aggregation heuristics: MIR kMaxAggregation=5
%   (linear_programming_constraint.cc:1924), CG via GetUnitRowLeftInverse (1654),
%   zero-half +/-1 (zero_half_cuts.h), clique via WeightedBronKerboschBitsetAlgorithm.
\textbf{One recipe behind many of them.} Despite the different names, many
generic cuts follow the same four-step pattern (\srcf{cuts.cc}).

\emph{(1) Aggregate.} Take a linear combination of LP rows (the same algebraic
operation behind the LP's reasons in \cref{sec:lp}) to obtain a single
constraint that is tight, or nearly tight, at the current fractional point. The
multipliers distinguish the families: Chvátal--Gomory reads them from a row of
$B^{-1}$, MIR combines only a handful of rows, and zero-half uses small
coefficients such as $\pm 1$.

\emph{(2) Rewrite.} Put the aggregate into a canonical form over non-negative
variables, complementing some variables ($x \mapsto \text{range}-x^c$) and
substituting implied bounds so that the rounding step is valid.

\emph{(3) Round.} Apply a \textbf{super-additive function}
$f:\mathbb{Z}\to\mathbb{Z}$ to every coefficient and to the right-hand side,
where super-additivity means $f(a)+f(b)\le f(a+b)$. This property is exactly
what keeps the rounded inequality valid for all integer points while allowing it
to cut off the current fractional point. Cover and knapsack cuts use a simple
function such as $f(x)=\lfloor x/d\rfloor$, while MIR uses
\[
  f(x)=s\lfloor x/d\rfloor+\max\!\bigl(0,(x \bmod d)-r\bigr).
\]

\emph{(4) Substitute back.} Translate the rounded inequality back into the
original variables.

Thus, a large part of the generic cut machinery reduces to an integer rounding
helper, \srcla{cuts.h}{341}{GetSuperAdditiveRoundingFunction}, wrapped in different
aggregation heuristics. This is also why the kept cuts remain exact: the
rounding function maps integers to integers, and the inequality that is added to
the LP is stored in integer arithmetic.

\textbf{Structural and scheduling cuts.} Because CP-SAT still \emph{knows} the
high-level constraints from \cref{sec:expansion}, it can add cuts that a
flat MIP model would not easily recover. For
\srcla{circuit.h}{46}{Circuit}/routing constraints, it can generate dynamic
\textbf{subtour-elimination} and capacity cuts. For
\srcla{disjunctive.h}{49}{NoOverlap} and
\srcla{cumulative.h}{48}{Cumulative}, it has a rich scheduling family:
\textbf{energetic} cuts, stating that the energy in a time window must fit
capacity times span, lifted by tasks that must partly overlap the window;
\textbf{completion-time} cuts, including the Smith-rule area bound
\[
  \sum_i p_i e_i \ge \frac12\left(\sum_i p_i^2 + \left(\sum_i p_i\right)^2\right)
\]
(here $p_i$ is the constant processing time and $e_i$ the completion-time
variable, so the inequality is linear in the $e_i$ with a constant right-hand
side); as well as time-table and precedence cuts. These families are an important part
of why CP-SAT is competitive on scheduling problems against dedicated solvers.

\textbf{Still exact, hence still compatible with learning.} The crucial property
is that cuts inherit the LP's exactness discipline from \cref{sec:lp}:
coefficients are kept in integer arithmetic (\code{IntegerValue}, with an
\code{int128} right-hand side), so every cut is a \emph{provably valid} integer
inequality, not a floating-point approximation. A cut is not itself a learned
clause; it is a valid linear row added to the relaxation
(\srcla{linear\_constraint\_manager.h}{154}{LinearConstraintManager}). But
because it is exact, it can participate cleanly in the LP's certified reasoning.
When the strengthened LP later derives a dual bound or a dual-ray conflict
(\cref{sec:lp}), the exact integer certificate behind that deduction may
use the cut, and \emph{that} reason is what enters conflict analysis and is
learned. Thus, cuts participate in learning indirectly, through the LP's exact
reasons, rather than as clauses on the trail.

\textbf{Management and cost.} Cuts are not free: each one enlarges the LP and
can slow subsequent solves. Generation is therefore concentrated at the root,
``level zero,'' where a tighter bound benefits the entire search below. The cut
pool is curated: candidates are filtered by \textbf{efficacy}, meaning how deeply
they cut the current point, and by \textbf{orthogonality}, meaning how much new
direction they add relative to existing cuts. The number of active cuts is capped
(\srcf{linear\_constraint\_manager.cc}), and constraints that remain unused or
basic for several solves can be dropped again. As with the LP itself, the
portfolio hedges how aggressively to invest: \code{max\_lp} leans into cuts,
whereas \code{no\_lp} forgoes them entirely.

% Source (VERIFIED): "Lp stats" column "AddedCuts" (stat_tables.cc:49); per-type
%   "Lp Cut" table from manager.type_to_num_cuts() (stat_tables.cc:282-298,395);
%   "Lp pool" column "Strengthened" = num_coeff_strengthening (stat_tables.cc:58,308).
\begin{platypuslog}[htb]
Cut activity is visible in two places: the \code{AddedCuts} column of the
\textsf{Lp stats} table gives the total, and a dedicated \textsf{Lp Cut} table
breaks it down by generator, such as \code{CG}, \code{MIR},
\code{ZERO\_HALF}, cover, clique, and the scheduling/routing families. The
coefficient strengthening from this section's opening appears one table over, as
the \code{Strengthened} column of the \textsf{Lp pool} table.
\end{platypuslog}
\subsection{Restarts: abandoning the tree without forgetting}\label{sec:restarts}
% Sources (VERIFIED):
%   ortools/sat/sat_parameters.proto -- RestartAlgorithm enum (NO_RESTART,
%     LUBY_RESTART, DL_MOVING_AVERAGE_RESTART, LBD_MOVING_AVERAGE_RESTART,
%     FIXED_RESTART); default_restart_algorithms =
%     "LUBY_RESTART,LBD_MOVING_AVERAGE_RESTART,DL_MOVING_AVERAGE_RESTART".
%   ortools/sat/restart.h / restart.cc -- RestartPolicy::ShouldRestart()/OnConflict();
%     LBD and decision-level moving averages.
%   ortools/sat/sat_solver.cc -- ComputeLbd(); on restart backtrack to level 0,
%     learned clauses are NOT cleared.
%   ortools/sat/sat_decision.h -- saved polarity (phase saving) and stable/
%     unstable phase (SetStablePhase); VSIDS-like variable activity.

A subtle weakness lurks in the loop of \cref{sec:loop}. Early decisions
are made when the solver knows least, yet they sit at the top of the tree and
constrain everything below. A bad first guess can trap search in a vast,
fruitless subtree for a long time. Worse, the branching heuristic's view of which
variables matter \emph{changes} as learning accumulates, but the decisions
already on the trail were made under an older, less-informed view.

\textbf{Restarts} are the cure. Periodically, the solver discards its current
decisions and jumps back to decision level 0, then starts descending again. The
critical detail, and the reason restarts help rather than merely spin, is what
is kept and what is discarded:

\begin{itemize}
  \item \textbf{Discarded:} the trail of decisions and the bounds they implied.
  The guesses are gone.
  \item \textbf{Kept:} \emph{everything the solver learned.} Learned nogood
  clauses, variable activity scores, and saved polarities all survive. In the
  code, a restart backtracks to level 0, but the learned-clause repository is not
  cleared (\srcla{sat\_solver.cc}{1734}{sat\_solver.cc}, the restart branch around
  line~1734).
\end{itemize}

A restart is therefore not a reset. It is a \emph{re-descent with accumulated
wisdom}. The solver returns to the root, but now with all the nogoods it has
learned. Propagation near the top of the tree is stronger than it was before,
and the branching heuristic, updated by variable activities, can make better
opening moves. Restarts also frequently make learned asserting clauses fire much
closer to the root, fixing facts that previously required deeper search.

\textbf{When to restart.} CP-SAT does not rely only on a fixed restart schedule;
it also watches the \emph{quality} of recent conflicts. The key quality metric is
\textbf{LBD}, Literal Block Distance: the number of distinct decision levels
appearing in a learned clause (\srcla{sat\_solver.cc}{1926}{ComputeLbd}).
A low-LBD clause ties together literals from few levels and tends to be highly
reusable. If recent LBDs trend worse than the running average, the current
search trajectory is degrading, and a restart becomes attractive.

CP-SAT cycles through a small set of restart strategies
(\srcla{sat\_parameters.proto}{267}{RestartAlgorithm}): a \textbf{Luby}
sequence, with intervals $1, 1, 2, 1, 1, 2, 4, \ldots$ and robust worst-case
properties; an \textbf{LBD moving-average} trigger; and a
\textbf{decision-level moving-average} trigger. The default configuration cycles
through all three (\code{default\_restart\_algorithms}), the aim being robustness
across instances that reward either aggressive or conservative restarting; other
workers may instead lean on a single algorithm.

\textbf{Phase saving and stability.} When a restart discards the trail, which
value should the solver try first for a variable it sees again? It remembers the
\textbf{last value that variable held}, its saved \emph{polarity}, and reuses it.
A restart therefore need not redo work from scratch: variables drift back toward
a partial assignment that was working, while early decisions are reconsidered
with better information.

CP-SAT modulates this behavior with a \textbf{stable/unstable} phase distinction
(\srcf{sat\_decision.h}). In \emph{stable} mode, paired with Luby restarts, it
restarts less aggressively and trusts saved phases, behaving more like a solver
trying to \emph{find} a solution. In \emph{unstable} mode, it restarts more often
and explores more broadly, behaving more like a solver trying to \emph{prove}
infeasibility. The solver alternates between the two rather than commit to either.

A restart lets CDCL escape bad early commitments while monotonically accumulating
the learned constraints that make each later descent cheaper. It changes
\emph{where} the solver searches without discarding anything it has proved.

% Source (VERIFIED): "Search stats" column "Restarts" = FormatCounter(counters.num_restarts)
%   at ortools/sat/stat_tables.cc:93. The response field "restarts:"
%   (cp_model_solver.cc:722) is one worker's count, not a portfolio sum (see the
%   portfolio box below); a 24-worker log showed restarts:0 in the summary while
%   quick_restart_no_lp restarted 23'060 times.
\begin{platypuslog}[htb]
The \code{Restarts} column of the \textsf{Search stats} table records how often
each worker jumped back to level~0. Aggressive configurations such as
\code{quick\_restart} often restart far more than steadier ones, although on some
models every worker restarts a similar number of times. This is one quantitative
glimpse of the portfolio's diversity.
\end{platypuslog}

\subsection{The portfolio: many configurations, one core}\label{sec:portfolio}
% Sources (VERIFIED):
%   ortools/sat/cp_model_search.cc -- builds the default subsolver portfolio
%     (default_lp, no_lp, max_lp, core, quick_restart, reduced_costs,
%     pseudo_costs, fixed, probing, objective_lb_search, objective_shaving,
%     lb_tree_search, ...), each a SatParameters variant of the same core.
%   ortools/sat/synchronization.h -- SharedClausesManager (shares learned
%     clauses, esp. short/binary) and SharedBoundsManager (shares improved
%     variable/objective bounds) between workers.
%   ortools/sat/cp_model_solver.cc -- worker setup and Synchronize().

CP-SAT's parallelism is the second half of the same idea. Given several workers
(threads), the natural expectation is ``divide the search tree among them.''
CP-SAT usually does something more robust: it runs a \textbf{portfolio}, where
each worker is a \emph{differently configured copy of the same search core}
attacking the \emph{whole} problem, and the workers cooperate by sharing what
they learn.

This is a \textbf{configuration portfolio, not an algorithm portfolio}. Every
worker runs the loop of \cref{sec:loop}; the workers differ mainly in
their parameters (\srcf{cp\_model\_search.cc} builds the default set). The exact
roster and the names below are specific to the pinned release and change between
versions; treat them as illustrative of the design, not as a stable contract. A
representative slice of the named full-thread subsolvers is:

\begin{itemize}
  \item \textbf{\code{default\_lp}}: the balanced default, combining propagation
  with a light LP relaxation that feeds bounds back into the integer trail.
  \item \textbf{\code{no\_lp}}: pure CP/SAT, without the LP relaxation; fast and
  lean when the LP does not pay off.
  \item \textbf{\code{max\_lp}}: a heavier LP relaxation for problems where
  strong linear bounds dominate.
  \item \textbf{\code{core}}: core-guided optimization, which attacks the
  objective through infeasible cores, possibly minimized heuristically, rather
  than only through linear bound-tightening; often effective on problems with
  many soft constraints. It is scheduled by default only for optimization
  problems whose objective has several terms, and dropped otherwise.
  \item \textbf{\code{quick\_restart}}: a configuration that restarts very
  aggressively, useful on problems that reward rapid re-descent.
  \item \textbf{\code{pseudo\_costs} / \code{reduced\_costs}}: configurations
  that branch using objective-derived estimates of which variables most affect
  cost, in a MIP-like style.
  \item \textbf{\code{fixed}}: a configuration that follows a user-provided
  fixed search order.
  \item \textbf{\code{probing}, \code{objective\_lb\_search},
  \code{objective\_shaving}, \code{lb\_tree\_search}}: configurations
  specializing in proving bounds rather than finding solutions.
\end{itemize}

The instance landscape is unpredictable: sometimes the LP is decisive, sometimes
core-guided search cracks the instance, and sometimes aggressive restarts win.
Running these configurations \emph{simultaneously} makes the wall-clock behavior
closer to the best configuration for that instance, without requiring you to know
in advance which one that is. This diversity is also one reason CP-SAT can see
superlinear speedups from additional cores: extra workers are not merely splitting
one fixed amount of work, but placing additional independent bets, any of which
may crack the instance.

% Sources (VERIFIED): incomplete (interleaved) subsolvers are built in
%   cp_model_solver.cc (SubSolver::INCOMPLETE; "feasibility_pump" :1223,
%   generator-named LNS/LS workers :1317; listed as "interleaved subsolver" :794;
%   LNS roster built :1879-2016 -- rins/rens, rnd/graph_*_lns, lb_relax_lns,
%   scheduling_*_lns, packing_*_lns).
%   use_feasibility_jump field 265 default=true (sat_parameters.proto:1292);
%   num_violation_ls field 244 (proto:1337); use_lns field 283 default=true
%   (proto:1483); NeighborhoodGenerator base class cp_model_lns.h:388.
\begin{platypuswarning}[htb][Incomplete search is part of the engine]
The portfolio configurations named so far are only the \emph{complete} (exhaustive) slice. Alongside them,
CP-SAT interleaves \textbf{incomplete} primal heuristics that trade completeness
for speed; the log calls them \emph{interleaved subsolvers}. They come in the two
classic flavors:
\begin{itemize}
  \item \textbf{First-solution heuristics} drive toward an initial feasible
  point. CP-SAT runs the \textbf{feasibility pump} and \textbf{feasibility jump}
  (\srcla{sat\_parameters.proto}{1292}{use\_feasibility\_jump}), a local
  search that repeatedly nudges an assignment to reduce total constraint
  violation.
  \item \textbf{Improvement heuristics} take an existing incumbent and make it
  better. The dominant family here is \textbf{Large Neighborhood Search}
  (\srcla{cp\_model\_lns.h}{388}{NeighborhoodGenerator}; \srcla{sat\_parameters.proto}{1483}{use\_lns}): fix most of the incumbent and let the full
  solver re-optimize a small neighborhood. CP-SAT ships many such neighborhoods,
  generic (\code{rins/rens}, random and graph-based) and structure-aware
  (scheduling- and packing-specific), complemented by constraint-based
  (violation) local search (\code{num\_violation\_ls}).
\end{itemize}
These never prove optimality, but a single good incumbent any of them publishes
through the shared-bounds channel tightens every other worker's objective cutoff
at once, so a worker that ``only'' finds or improves solutions also accelerates
the optimality proof run by the exact workers. The incomplete heuristics are
therefore a fundamental source of CP-SAT's practical strength, not a sideshow
bolted onto the exact core.

For the curious reader, the constraint-based local-search worker is described in
a paper by the developers \cite{davies2024violationls}; and \cite{berthold2025primal}
is a compact, readable book covering both families above, including Large
Neighborhood Search, for integer programming in general.
\end{platypuswarning}

\textbf{Cooperation is what makes it more than a race.} The workers are not
isolated. Two shared channels (\srcf{synchronization.h}) let a discovery by one
worker accelerate the others:

\begin{itemize}
  \item \textbf{Shared clauses}
  (\srcla{synchronization.h}{853}{SharedClausesManager}): learned nogoods,
  especially short, low-LBD clauses and binary clauses, are broadcast between
  workers. A dead end that \code{core} discovers can become a permanent
  constraint that \code{default\_lp} and \code{no\_lp} also benefit from. This
  turns the portfolio into a \emph{collective} learner: the global pool of useful
  nogoods can grow faster than any single worker could produce.
  \item \textbf{Shared bounds}
  (\srcla{synchronization.h}{661}{SharedBoundsManager}): when one worker tightens
  a variable's domain or improves the objective bound, it publishes the
  improvement; the others import it and propagate it. In optimization, the moment
  any worker finds a better incumbent, every worker's objective cutoff
  $\code{obj} \le C-1$ tightens as well. Good solutions therefore propagate
  across the portfolio and accelerate the optimality proof elsewhere.
\end{itemize}

Parallel CP-SAT is therefore best understood as \textbf{one search core,
diversified into many parameterizations, all solving the full problem and
pooling their learned clauses and bounds.} The diversity covers uncertainty
about which configuration suits the instance; the sharing ensures that a
worker's hard-won insight is not wasted on the others.

% Sources (VERIFIED against sat_parameters.proto):
%   num_workers field 206, default 0 (=auto): sat_parameters.proto:707
%     ("num_workers is the preferred name", doc :695-706).
%   subsolvers (repeated string) field 207: proto:742 -- names the full-thread
%     portfolio configs; doc :717-741 lists default_lp/no_lp/max_lp/quick_restart/
%     pseudo_costs/reduced_costs/probing/lb_tree_search/fixed and notes the order
%     matters (only the first num_full_subsolvers are scheduled).
%   extra_subsolvers (repeated string) field 219: proto:746 ("add more workers
%     types ... added at the beginning of the list").
%   ignore_subsolvers (repeated string) field 209: proto:757 -- glob patterns
%     ('*','?') to drop configs; doc :751-756.
%   linearization_level field 90, default 1: proto:1652 (0 = no LP, 2 = heavy LP;
%     this is exactly what no_lp / max_lp set, per the subsolvers doc :721-725).
\begin{platypusparam}[htb]
You are not stuck with the default mix. \srcla{sat\_parameters.proto}{742}{subsolvers}, a list of named configurations, lets you state
\emph{exactly} which full-thread workers run. \code{extra\_subsolvers} appends
your own, and \code{ignore\_subsolvers} drops configurations by glob pattern
(e.g., \code{"max\_lp"} to suppress the heavy LP on a model where it never pays).
\srcla{sat\_parameters.proto}{707}{num\_workers} sets how many workers run
at once. Pin it to \code{1} with a single subsolver to get a deterministic,
single-strategy search, the lever that makes the ``restrict which workers run''
caveat concrete. Because the named configurations differ largely in how
hard they lean on the LP, \srcla{sat\_parameters.proto}{1652}{linearization\_level} (\code{0} = none, like \code{no\_lp};
\code{2} = heavy, like \code{max\_lp}) dials that axis directly on whichever
workers you keep.
\end{platypusparam}

% Source (VERIFIED): every stat table is keyed by FormatName(name) per subsolver
%   (ortools/sat/stat_tables.cc, e.g. search_table_/lp_table_ push one row per worker);
%   progress lines carry the subsolver via ExtractSubSolverName()/RegisterSolutionFound()
%   in cp_model_solver_logging.cc.
% Source (VERIFIED): the closing response summary is NOT a portfolio total. Its
%   counters are written by a single statistics postprocessor run on ONE model:
%   SharedResponseManager::FillSolveStatsInResponse(model,...) (synchronization.cc:271)
%   reads that model's SatSolver (num_branches/num_failures/num_propagations/
%   num_restarts) and IntegerTrail (num_enqueues) -- cp_model_solver.cc:2317-2348;
%   the summary block just prints that one CpSolverResponse (cp_model_solver.cc:710-723).
% Empirical (xl_20x6_120parties.log, 24 workers): summary read conflicts:0
%   branches:1072 restarts:0, while the 'core' Search stats row recorded
%   921'777 conflicts / 10'142'481 branches / 530M BoolPropag / 100M IntegerPropag.
\begin{platypuslog}[htb]
The portfolio is why every end-of-run table carries \emph{one row per worker}:
the \textsf{Search stats}, \textsf{Lp stats}, and timing tables list
\code{default\_lp}, \code{no\_lp}, \code{max\_lp}, \code{core}, and the rest side
by side. Each \code{\#1}/\code{\#Bound} progress line is also tagged with the
subsolver that produced it, so the log shows which configuration actually
cracked the instance.

The same fact makes the closing \textsf{CpSolverResponse summary} easy to
misread. Its headline counters --- \code{booleans:}, \code{conflicts:},
\code{branches:}, \code{propagations:}, \code{integer\_propagations:},
\code{restarts:} --- are \emph{not} portfolio totals. Each is copied from the one
worker whose model carries the returned solution, which is frequently a
first-solution or LNS helper that barely searched. In one real 24-worker run the
summary read \code{conflicts: 0} and \code{branches: 1072} while the \code{core}
worker's \textsf{Search stats} row recorded about 920{,}000 conflicts and over
ten million branches. To see how hard the search actually worked, read the
per-worker table, not the summary.
\end{platypuslog}
